\section{Validering av provider-arkitekturen}
\label{sec:provider-validering}

For å undersøke om nye videomotorer kan legges til uten å endre
kjernelogikken, ble to tilleggsprovidere implementert i separate
feature-branches etter at AntMedia og NHN var ferdig.
\texttt{MockVideoProvider} var en stub uten eksterne avhengigheter.
\texttt{DailyProvider} var en reell integrasjon mot Daily.co sitt
REST-API. Ingen av disse var opprinnelige HDO-krav, men ble brukt
for å validere utvidbarheten i arkitekturen.

Endringsomfanget ble målt ved å sammenligne hver branch mot
\texttt{develop} med \texttt{git diff --stat},
\texttt{git diff --name-status} og \texttt{git diff --shortstat}.
Alle tre branchene hadde samme merge-base (\texttt{5dd99b5}).
Tabell~\ref{tab:provider-validering} viser resultatene.

\begin{table}[h]
\centering
\begin{tabular}{lrrrr}
\toprule
\textbf{Branch} & \textbf{Filer} & \textbf{Linjer +} &
\textbf{Linjer -} & \textbf{Kjerneendringer} \\
\midrule
\texttt{feature/mock-provider}       & 6  & 249  & 2 & 0 \\
\texttt{feature/daily-provider}      & 12 & 779  & 7 & 0 \\
\texttt{feature/provider-validation} & 14 & 1024 & 8 & 0 \\
\bottomrule
\end{tabular}
\caption{Endringsomfang per branch målt mot \texttt{develop}.
Kjerneendringer refererer til endringer i \texttt{RoomsController},
\texttt{IVideoProvider} og offentlige modeller.}
\label{tab:provider-validering}
\end{table}

I samtlige sammenligninger viste \texttt{git diff --name-status}
ingen endringer i \texttt{RoomsController}, \texttt{IVideoProvider}
eller de offentlige modellene. Begge providerene kunne dermed legges
til uten endringer i kontrollerlaget eller API-kontrakten.
Endringene var begrenset til nye provider-klasser, konfigurasjon,
dependency injection-registrering og tester.

Alle tre branchene ble bygget og testet lokalt med
\texttt{dotnet build} og \texttt{dotnet test} på macOS. Bygg og
tester passerte i alle tre branchene. Den samlede
\texttt{feature/provider-validation}-branchen ble i tillegg kjørt
på VM med Docker, der frontend ble brukt til å bekrefte at nye
providere dukket opp og fungerte mot API-et. E2E-testskriptet ble
ikke kjørt mot tilleggsproviderene — målet var å vise at
arkitekturen støtter utvidelse, ikke å teste providerenes
funksjonalitet i detalj.